\documentclass[10pt,letterpaper]{article}

\usepackage[letterpaper,margin=0.8in]{geometry}
\usepackage[T1]{fontenc}
\usepackage[utf8]{inputenc}
\usepackage{times}
\usepackage{url}
\usepackage[hidelinks]{hyperref}
\usepackage{cite}

\setlength{\parindent}{0pt}
\setlength{\parskip}{0.35em}
\renewcommand{\refname}{REFERENCES}
\sloppy

\makeatletter
\renewcommand{\maketitle}{%
  \begin{center}
    {\Large\bfseries \@title\par}
    \vspace{0.55em}
    {\normalsize \@author\par}
  \end{center}
  \vspace{-0.7em}
}
\renewcommand\section{\@startsection{section}{1}{\z@}{1.0em}{0.45em}{\normalfont\normalsize\bfseries}}
\renewcommand\subsection{\@startsection{subsection}{2}{\z@}{0.8em}{0.35em}{\normalfont\normalsize\bfseries}}
\makeatother


% IEEE-like author block commands while keeping a one-column article format.
\newcommand{\IEEEauthorblockN}[1]{\textbf{#1}}
\newcommand{\IEEEauthorblockA}[1]{#1}

\title{Context-Aware Alert Prioritization for Security Information and Event Management Systems}
\author{%
    \IEEEauthorblockN{Anawin Suwannapha}%
    , \IEEEauthorblockN{Worathun Termjitcharoen}\\
    \IEEEauthorblockA{\textit{School of Information, Computer and Communication Technology (ICT)}}\\
    \IEEEauthorblockA{\textit{Sirindhorn International Institute of Technology, Thammasat University, Thailand}}\\
    \IEEEauthorblockA{6622780573@g.siit.tu.ac.th, 6622790168@g.siit.tu.ac.th}
}
\date{}
\begin{document}

\maketitle

\section{INTRODUCTION}

Security Information and Event Management (SIEM) systems are widely used in modern Security Operations Centers (SOCs) to collect, correlate, and analyze security events, but large alert volumes can overwhelm analysts and delay the investigation of high-risk incidents \cite{ref14}. A major limitation is that alerts are often ranked mainly by static rule severity or identifiers, which may not reflect the real operational risk of an event when factors such as asset criticality, user behavior, vulnerability exposure, timing, and surrounding attack patterns are considered. Prior research has therefore explored context-aware alert classification, adaptive SIEM correlation, contextual log filtering, reinforcement learning-based prioritization, and explainable graph-based triage to reduce alert fatigue and help analysts focus on alerts with greater security impact \cite{ref1,ref5,ref6,ref7,ref8}.

Wazuh provides a suitable base for this type of research because it already supports log analysis, endpoint monitoring, rule-based detection, vulnerability detection, threat intelligence integration, MITRE ATT\&CK-aligned evaluation, and active response workflows. Recent Wazuh-focused studies have explored real-time threat feed integration \cite{ref3}, MITRE ATT\&CK evaluation \cite{ref12}, and machine learning-based improvements to threat detection \cite{ref13}. However, many of these approaches focus mainly on detection accuracy, threat intelligence enrichment, or response automation. Less attention has been given to an additional Wazuh-based prioritization layer that uses created context to decide which alerts should be investigated first.

This project proposes a context-aware alert prioritization framework built on top of Wazuh. The goal is not to replace Wazuh's existing detection rules or build a new SIEM platform, but to design an additional prioritization layer that enriches Wazuh alerts with technical security evidence and real-world operational context. The proposed framework will use alert and log information together with real-life context to create context such as the likely target of attack, possible blast radius, estimated attacker sophistication, duration of attack, unusual event timing, recent attack patterns, alert frequency, asset criticality, vulnerability exposure, threat intelligence reputation, MITRE ATT\&CK tactic, and historical alert behavior. This context will then be used to support a hybrid risk scoring model that combines Wazuh rule severity with contextual risk factors to produce an interpretable priority score for SOC analysts. Related risk-based prioritization work in SIEM environments supports the need for transparent scoring models that remain understandable to human analysts \cite{ref9}, while broader SIEM alert management research shows the value of combining multiple alert features during prioritization \cite{ref10}.

In addition to hybrid scoring, this project explores graph-based context link analysis. Alerts and related entities can be represented as an alert graph, where nodes represent alerts, assets, users, vulnerabilities, attack techniques, threat indicators, or other relevant entities, and edges represent contextual links created from the extracted context. A graph neural network can then be applied over this alert graph so that risk is propagated through connected alerts and entities. This allows low-severity alerts linked to high-risk neighbors to be promoted when they are part of a wider suspicious pattern. This direction follows recent Wazuh graph-based research on real-time attack path prediction using event correlation \cite{ref4} and Wazuh EDR telemetry research that applies graph neural networks to model relationships among security-relevant events and endpoint behaviors \cite{ref11}.

The expected contribution of this work is a Wazuh-based context-aware alert prioritization framework that supports SOC triage by reducing repeated low-value alerts and promoting high-risk alerts that may otherwise be hidden in large alert volumes. The project contributes a practical context creation method for transforming Wazuh alerts and related logs into useful prioritization features, a hybrid risk scoring model that remains interpretable for analysts, and a graph-based link analysis approach for associating related alerts into broader incidents.

\section{PROJECT SCOPE}

The scope of this project includes the following key areas:

\begin{enumerate}
\item \textbf{Wazuh-Based Alert Prioritization:} Designing an additional prioritization layer on top of Wazuh, rather than replacing Wazuh's existing detection rules or building a new SIEM platform.

\item \textbf{Alert Log Enrichment:} Enhancing Wazuh alerts with information such as rule level, MITRE ATT\&CK tactic, alert frequency, asset criticality, vulnerability exposure, threat intelligence reputation, and historical alert behavior.

\item \textbf{Context Creation:} The framework computes context directly from Wazuh alert history using two complementary extraction methods. First, a sliding-window behavioral profile aggregates each entity's, such as source IP, user, or host, most recent alerts into features \cite{ref15}. Second, structural context such as blast radius and attack-target centrality is derived from the alert-entity graph. These generated contexts are then combined with looked-up and real-life operational context, such as asset criticality, vulnerability exposure, threat-intelligence reputation, business role of the affected asset, and expected user or system behavior \cite{ref14}.

\item \textbf{Graph Neural Network:} Applying a graph neural network over an alert graph whose nodes are alerts and their associated entities, and whose edges are the contextual links created in item 3. Message passing propagates risk across connected alerts, so that low-severity alerts linked to high-risk neighbors are promoted, following recent graph-based work on Wazuh telemetry \cite{ref4,ref11}.

\item \textbf{Analyst-Focused Output:} Producing a prioritized alert list that supports SOC triage and helps reduce analyst workload caused by repeated or low-value alerts.
\end{enumerate}

\section{CONTRIBUTIONS}

This work makes the following key contributions:

\begin{enumerate}
\item \textbf{Context-Aware Wazuh Alert Prioritization Framework:} A hybrid framework that enriches Wazuh alerts with operational context such as MITRE ATT\&CK mapping, asset criticality, vulnerability exposure, threat intelligence, alert frequency, and historical behavior.

\item \textbf{Hybrid Risk Scoring Model:} A transparent scoring method that combines Wazuh rule severity with contextual risk factors to produce an interpretable priority score for SOC analysts.

\item \textbf{Context Link Analysis:} Using the created contexts as linking keys to associate related alerts into incidents, and applying link prediction over the resulting graph to uncover relationships that individual detection rules do not express.

\item \textbf{Alert Fatigue Reduction:} A prioritization mechanism designed to reduce repeated low-value alerts and promote high-risk alerts that may otherwise be hidden in large alert volumes.

\item \textbf{Wazuh-Based Evaluation:} An evaluation approach comparing default Wazuh severity ranking against the proposed context-aware ranking, with possible use of MITRE ATT\&CK-aligned attack simulations.
\end{enumerate}

\begin{thebibliography}{15}
\scriptsize
\setlength{\itemsep}{0pt}
\setlength{\parskip}{0pt}
\setlength{\parsep}{0pt}

\bibitem{ref1} J. Roy and E. A. Balde, ``Toward Context-Aware Alert Classification in Security Operations Centers Using LLMs,'' 2026 IEEE 5th International Conference on AI in Cybersecurity (ICAIC), Houston, TX, USA, 2026, pp. 1--10, doi: 10.1109/ICAIC67076.2026.11395735.

\bibitem{ref2} A. Hassan, A. Rauf, N. Shafqat, R. Latif, and H. Khan, ``ZenGuard a machine learning based zero trust framework for context aware threat mitigation using SIEM SOAR and UEBA,'' Scientific Reports, vol. 15, 2025, doi: 10.1038/s41598-025-20998-4.

\bibitem{ref3} A. Al Siam, M. M. Hassan, A. K. M. Masum, and T. Bhuiyan, ``Automating Malware Detection and Response via Real-Time Threat Feed Integration with Wazuh SIEM,'' 2025 IEEE 2nd International Conference on Computing, Applications and Systems (COMPAS), Kushtia, Bangladesh, 2025, pp. 1--6, doi: 10.1109/COMPAS67506.2025.11381876.

\bibitem{ref4} S. Chaurasia, A. Singh, S. K. Maurya, D. P. Bajpai, and S. Yadav, ``Graph Neural Threat Intelligence: Real-Time Attack Path Prediction Using Event Correlation from Wazuh SIEM,'' 2026 International Conference on Intelligent Processing, Hardware, Electronics, and Radio Systems (CIPHER), Jalandhar, India, 2026, pp. 1--5, doi: 10.1109/CIPHER70417.2026.11523812.

\bibitem{ref5} X. Wang, X. Yang, X. Liang, X. Zhang, W. Zhang, and X. Gong, ``Combating alert fatigue with AlertPro: Context-aware alert prioritization using reinforcement learning for multi-step attack detection,'' Computers \& Security, vol. 137, 2024, Art. no. 103583, doi: 10.1016/j.cose.2023.103583.

\bibitem{ref6} G. Suarez-Tangil, E. Palomar, A. Ribagorda, and I. Sanz, ``Providing SIEM systems with self-adaptation,'' Information Fusion, vol. 21, pp. 145--158, 2015, doi: 10.1016/j.inffus.2013.04.009.

\bibitem{ref7} M. Cinque, R. Della Corte, and A. Pecchia, ``Contextual filtering and prioritization of computer application logs for security situational awareness,'' Future Generation Computer Systems, vol. 111, pp. 668--680, 2020, doi: 10.1016/j.future.2019.09.005.

\bibitem{ref8} A. Malach, P. N. Wudali, S. Momiyama, J. Furukawa, T. Araki, Y. Elovici, and A. Shabtai, ``CyberShapley: Explanation, prioritization, and triage of cybersecurity alerts using informative graph representation,'' Computers \& Security, vol. 150, 2025, Art. no. 104270, doi: 10.1016/j.cose.2024.104270.

\bibitem{ref9} Z. Shamsi, N. V. S. Manohar, P. Likhitha, M. J. S. S. S. Jahnavi, and K. E. Rao, ``A Human Centric Risk-Based Security Alert Prioritization Using Splunk,'' 2026 IEEE Guwahati Subsection Conference (GCON), Guwahati, India, 2026, pp. 1--7, doi: 10.1109/GCON69192.2026.11649152.

\bibitem{ref10} M. Khayat, E. Barka, M. A. Serhani, F. Sallabi, K. Shuaib, and H. M. Khater, ``Advanced Techniques for Alert Management in Security Information and Event Management Systems With Ensembled Deep Learning, Hybrid Optimization, and Multi-Feature Extraction,'' IEEE Open Journal of the Communications Society, vol. 6, pp. 7349--7368, 2025, doi: 10.1109/OJCOMS.2025.3603000.

\bibitem{ref11} A. Pomsathit, ``Adaptive Detection of Advanced Persistent Threats (APT) With Graph Neural Networks and Rehearsal-Based Continual Learning on Wazuh EDR Telemetry,'' IEEE Access, vol. 13, pp. 212973--212982, 2025, doi: 10.1109/ACCESS.2025.3639270.

\bibitem{ref12} A. M. Winkler and P. Sharma, ``Proactive threat detection in enterprise systems using Wazuh: A MITRE ATT\&CK Evaluation,'' Computers \& Security, vol. 159, 2025, Art. no. 104702, doi: 10.1016/j.cose.2025.104702.

\bibitem{ref13} S. A. Chamkar, M. Zaydi, Y. Maleh, and N. Gherabi, ``Improving Threat Detection in Wazuh Using Machine Learning Techniques,'' Journal of Cybersecurity and Privacy, vol. 5, no. 2, Art. no. 34, 2025, doi: 10.3390/jcp5020034.

\bibitem{ref14} S. Ndichu, T. Ban, S. Ozawa, T. Takahashi, and D. Inoue, ``AI-Driven Security Alert Screening and Alert Fatigue Mitigation in Security Operations Centers: A Survey,'' arXiv:2605.08316 [cs.CR], May 2026, doi: 10.48550/arXiv.2605.08316.

\bibitem{ref15} T. van Ede et al., ``DEEPCASE: Semi-Supervised Contextual Analysis of Security Events,'' 2022 IEEE Symposium on Security and Privacy (SP), San Francisco, CA, USA, 2022, pp. 522--539, doi: 10.1109/SP46214.2022.9833671.

\end{thebibliography}

\end{document}